\SourceLeaf{460}
\begin{center}\Large\bfseries BEMERKUNGEN ZUM ACHTEN BANDE.\end{center}

Der vorliegende achte Band von \textsc{Gauss}' Werken gibt Ergänzungen zu den drei ersten Bänden und zum vierten Bande mit Ausnahme des geodätischen Theils. Er enthält:

unter den Abtheilungen \textit{Arithmetik, Analysis} und \textit{Numerisches Rechnen} eine Reihe einzelner Notizen und Aufsätze, die sich nachträglich in \textsc{Gauss}' Nachlass vorgefunden haben,

aus der \textit{Wahrscheinlichkeitsrechnung} hauptsächlich eine Anzahl Belege über die Entdeckung der Methode der kleinsten Quadrate;

aus der \textit{Geometrie} Notizen und Briefe, die sich auf die Grundlagen der Geometrie, die Geometria situs, Lehrsätze und Aufgaben aus der elementaren Geometrie, die Verwendung complexer Grössen für die Geometrie und die Theorie der krummen Flächen beziehen.

Die Bearbeitung der Abtheilungen \textit{Arithmetik und Algebra} und \textit{Analysis und Functionentheorie} rührt von Herrn \textsc{Fricke} in Braunschweig, die der Abtheilungen \textit{Numerisches Rechnen} und \textit{Wahrscheinlichkeitsrechnung} von den Herren \textsc{Börsch} und \textsc{Krüger} in Potsdam und die der \textit{geometrischen} Abtheilungen von Herrn \textsc{Stäckel} in Kiel her; die Notiz \textit{Allgemeinste Auflösung des Problems der Abwickelung der Flächen} (Seite 447 bis 449) ist von Herrn \textsc{Weingarten} in Charlottenburg bearbeitet worden.

Die allgemeine Redaction wurde von Herrn \textsc{Brendel} in Göttingen besorgt.

\textsc{Schering}, welcher dem vierten Bande einen Ergänzungsband hinzuzufügen beabsichtigte, hat eine Reihe von Vorarbeiten hinterlassen, hauptsächlich Abschriften aus dem Nachlass und dem Briefwechsel, welche bei der Herausgabe benutzt werden konnten.\hfill\mbox{\textsc{F. Klein.}}
